% =========================================================================
% PLANTILLA MAESTRA: APUNTES TÉCNICOS MONOCROMO (1 COLUMNA COMPLETA)
% =========================================================================
% Optimizada para impresión física láser monocromática (HP LaserJet P1006 / Tóner).
%
% Principios de Diseño Editorial:
% 1. Formato a 1 Columna Continua: Ancho útil completo (162 mm) con compensación
%    segura para anillado o perforado (bindingoffset 6 mm).
% 2. Sin Encabezado Superior: Cero running headers y cero reglas superiores para
%    maximizar la superficie útil vertical de lectura y evitar distracciones visuales.
% 3. Sin Espacio de Margen Lateral: Se eliminan los márgenes exteriores de notas,
%    priorizando el despliegue estructurado de tablas, fórmulas y definiciones.
% 4. Cero Solapamiento en Cajas y Tablas: Títulos integrados dentro del flujo
%    superior interno en tcolorbox y tablas con auto-ajuste de texto en celdas.
% 5. Tipografía Láser de Alto Rendimiento: Bitstream Charter a 10pt/13pt para el
%    cuerpo, TeX Gyre Heros para títulos sans-serif y Courier para código.
% 6. Alto Contraste Monocromo: 100% negro puro (#000000) sobre blanco. Ahorro de tóner
%    y cero degradado sucio en impresoras láser de 600 DPI.
% =========================================================================

\documentclass[10pt,a4paper,twoside]{extarticle}
\usepackage[utf8]{inputenc}
\usepackage[T1]{fontenc}

% --- Tipografía Láser Monocromo ---
\usepackage{charter}                         % Bitstream Charter: Excelente visibilidad y contraste
\usepackage{tgheros}                         % TeX Gyre Heros (Helvetica) para títulos sans
\usepackage{courier}                         % Courier para parámetros y código monoespaciado
\usepackage[spanish,es-tabla,es-noquoting]{babel}
\usepackage{microtype}                       % Ajuste óptico de justificación

% --- Geometría: Sin Margen Lateral y Sin Encabezado Superior ---
\usepackage{geometry}
\geometry{
    top=1.8cm,
    bottom=1.8cm,
    inner=2.0cm,
    outer=1.6cm,
    bindingoffset=6mm,                       % Margen seguro para perforado o anillado
    textwidth=162mm,                         % Ancho completo utilizable en 1 sola columna
    marginparwidth=0mm,                      % Margen lateral exterior completamente suprimido
    marginparsep=0mm,                        % Separador de margen lateral suprimido
    headheight=0pt,                          % Sin altura de encabezado
    headsep=0pt,                             % Sin separación de encabezado
    footskip=9mm                             % Pie de página compacto y sobrio
}

% --- Interlineado y Párrafos ---
\linespread{1.14}                            % Lectura cómoda sin dispersión
\setlength{\parindent}{0pt}
\setlength{\parskip}{4.5pt}
\raggedbottom

% --- Paquetes Complementarios ---
\usepackage{amsmath,amssymb}
\usepackage{booktabs}                        % Tablas profesionales sin líneas verticales
\usepackage{tabularx}                        % Tablas con auto-ajuste de ancho y envoltura
\usepackage{enumitem}
\usepackage{fancyhdr}
\usepackage{titlesec}
\usepackage{listings}
\usepackage[many]{tcolorbox}
\usepackage{hyperref}

% --- Encabezados y Pies de Página (Sin Encabezado Superior) ---
\pagestyle{fancy}
\fancyhf{}
\fancyhead{}                                 % Encabezado superior estrictamente vacío
\renewcommand{\headrulewidth}{0pt}           % Regla superior suprimida

\fancyfoot[LE,RO]{\fontsize{8.5pt}{10.5pt}\selectfont\sffamily\bfseries \thepage}
\fancyfoot[RE,LO]{\fontsize{7.5pt}{9.5pt}\selectfont\sffamily\color{black!70} [Materia / Asignatura] $\cdot$ Compendio Técnico}
\renewcommand{\footrulewidth}{0.3pt}

\fancypagestyle{firstpage}{
    \fancyhf{}
    \fancyfoot[LE,RO]{\fontsize{8.5pt}{10.5pt}\selectfont\sffamily\bfseries \thepage}
    \fancyfoot[RE,LO]{\fontsize{7.5pt}{9.5pt}\selectfont\sffamily\color{black!70} [Materia / Asignatura] $\cdot$ Edición Monocromo}
    \renewcommand{\headrulewidth}{0pt}
    \renewcommand{\footrulewidth}{0.3pt}
}

% --- Formato de Títulos (Jerárquicos y Sin Solapamiento) ---
\titleformat{\section}
  {\fontsize{12.5pt}{15.5pt}\bfseries\sffamily\color{black}}
  {\thesection.}{0.5em}{}
  [\vspace{2pt}\titlerule]
\titlespacing*{\section}{0pt}{11pt}{4pt}

\titleformat{\subsection}
  {\fontsize{10.5pt}{13.5pt}\bfseries\sffamily\color{black!90}}
  {\thesubsection}{0.5em}{}
\titlespacing*{\subsection}{0pt}{8pt}{3pt}

\titleformat{\subsubsection}
  {\fontsize{9.5pt}{12pt}\bfseries\sffamily\color{black!80}}
  {\thesubsubsection}{0.5em}{}
\titlespacing*{\subsubsection}{0pt}{6pt}{2pt}

% --- Listas Compactas ---
\setlist[itemize]{leftmargin=1.4em, itemsep=2pt, topsep=2pt, label=$\triangleright$}
\setlist[enumerate]{leftmargin=1.4em, itemsep=2pt, topsep=2pt}

% --- Cajas Técnicas con Protección de Borde (Cero Solapamiento de Título) ---
% 1. Definición Canónica: Barra lateral negra pura, título integrado en flujo superior interno
\newtcolorbox{definicion}[1][]{
    enhanced,
    breakable,
    colback=white,
    colframe=black,
    arc=0mm,
    boxrule=0pt,
    leftrule=2.4pt,
    top=5pt, bottom=5pt, left=9pt, right=8pt,
    before skip=6pt, after skip=6pt,
    fonttitle=\fontsize{9.5pt}{12pt}\selectfont\sffamily\bfseries,
    coltitle=black,
    title={Definición: #1},
    attach title to upper={\par\vspace{2.5pt}},
    fontupper=\fontsize{9.3pt}{12.5pt}\selectfont
}

% 2. Alerta de Examen / Punto Crítico: Borde continuo con separador interno seguro
\newtcolorbox{alertaparcial}[1][]{
    enhanced,
    breakable,
    colback=white,
    colframe=black!85,
    arc=0.8mm,
    boxrule=0.6pt,
    top=6pt, bottom=6pt, left=9pt, right=9pt,
    before skip=7pt, after skip=7pt,
    fonttitle=\fontsize{9.2pt}{11.5pt}\selectfont\sffamily\bfseries,
    coltitle=black,
    title={$\blacktriangleright$\ Punto Crítico / Criterio de Examen: #1},
    attach title to upper={\par\vspace{3pt}\hrule height 0.4pt\vspace{4pt}},
    fontupper=\fontsize{9.2pt}{12.5pt}\selectfont
}

% 3. Caja de Fórmulas y Demostraciones Matemáticas
\newtcolorbox{formulabox}[1][]{
    enhanced,
    colback=white,
    colframe=black!80,
    arc=0.6mm,
    boxrule=0.5pt,
    top=5pt, bottom=5pt, left=8pt, right=8pt,
    before skip=5pt, after skip=5pt,
    fontupper=\fontsize{9.2pt}{12.5pt}\selectfont
}

% --- Configuración de Código Fuente ---
\lstdefinestyle{monocodestyle}{
    basicstyle=\ttfamily\fontsize{8pt}{10pt}\selectfont\color{black},
    keywordstyle=\bfseries\color{black},
    commentstyle=\itshape\color{black!70},
    stringstyle=\color{black!90},
    showstringspaces=false,
    showspaces=false,
    showtabs=false,
    frame=single,
    framerule=0.4pt,
    rulecolor=\color{black!70},
    backgroundcolor=\color{white},
    breaklines=true,
    tabsize=2,
    aboveskip=5pt,
    belowskip=5pt,
    xleftmargin=6pt,
    xrightmargin=6pt
}

% Configuración de Tablas (Separación Vertical de Celdas)
\setlength{\extrarowheight}{3pt}
\renewcommand{\arraystretch}{1.2}

\hypersetup{
    colorlinks=false,
    pdfborder={0 0 0}
}

% =========================================================================
% INICIO DEL DOCUMENTO
% =========================================================================
\begin{document}
\thispagestyle{firstpage}

% --- Encabezado Superior de Portada Compacta ---
\noindent
{\fontsize{16pt}{20pt}\selectfont\sffamily\bfseries [TÍTULO PRINCIPAL DE LA ASIGNATURA]}\\[3pt]
{\fontsize{10.5pt}{13.5pt}\selectfont\sffamily\color{black!80} [Subtítulo o Módulo Temático -- Compendio Integral de Cátedra]}\\[2pt]
{\fontsize{8.5pt}{11pt}\selectfont\color{black!65} Facultad de Ingeniería $\cdot$ Universidad Nacional de Jujuy (UNJu) $\cdot$ Edición Monocromo Láser}\\[4pt]
\rule{\textwidth}{0.8pt}
\vspace{2mm}

% =========================================================================
% SECCIÓN DE EJEMPLO
% =========================================================================
\section{Introducción y Conceptos Fundamentales}

\subsection{1.1. Marco Teórico Principal}
Texto introductorio del tema desarrollado. El formato a 1 columna continua con tipografía \textbf{Bitstream Charter} a 10pt garantiza una densidad de lectura óptima en papel, permitiendo que el texto fluya con total naturalidad sin recortes laterales.

\begin{definicion}[Concepto Primario]
Explicación formal y canónica del concepto según la doctrina de la cátedra universitaria:
\begin{itemize}
    \item \textbf{Elemento A:} Descripción técnica detallada del primer componente.
    \item \textbf{Elemento B:} Descripción técnica del segundo componente.
    \item \textbf{Elemento C:} Justificación e impacto dentro del sistema integral.
\end{itemize}
\end{definicion}

\subsection{1.2. Tabla Comparativa de Modelos}
Las tablas utilizan el entorno \texttt{tabularx} para distribuir equitativamente el ancho útil de la página sin desbordamientos de texto:

\begin{table}[h!]
\centering
\small
\begin{tabularx}{\textwidth}{@{} >{\bfseries\raggedright\arraybackslash}p{3.2cm} >{\raggedright\arraybackslash}p{5.5cm} >{\raggedright\arraybackslash}X @{}}
\toprule
Criterio Evaluado & Enfoque Convencional & Enfoque Avanzado / Cátedra \\
\midrule
\textbf{Objetivo Primario} & Descripción del mecanismo base y sus limitaciones operativas. & Mecanismo optimizado de alta disponibilidad y tolerancia a fallos. \\
\addlinespace
\textbf{Tiempo de Respuesta} & Latencia estándar dependiente de conmutación electromecánica. & Conmutación instantánea ($0\text{ ms}$) mediante doble conversión continua. \\
\addlinespace
\textbf{Resiliencia} & Dependencia de intervención manual para restauración de servicio. & Recuperación automatizada mediante tableros de transferencia y redundancia. \\
\bottomrule
\end{tabularx}
\end{table}

\begin{alertaparcial}[Advertencia de Evaluación Parcial]
Recordar que en exámenes teóricos la cátedra suele evaluar la distinción precisa entre mecanismos preventivos y correctivos, haciendo énfasis en que las medidas activas reducen la probabilidad del siniestro, mientras que las pasivas mitigan el impacto.
\end{alertaparcial}

\subsection{1.3. Expresiones Matemáticas y Métricas}
Las fórmulas cuantitativas se presentan enmarcadas con líneas finas de alto contraste:

\begin{formulabox}
\textbf{Cálculo de Pérdida Anual Esperada (ALE):}
$$ALE = SLE \times ARO$$
Donde $SLE$ representa la pérdida ante un evento individual ($SLE = AV \times EF$) y $ARO$ es la frecuencia anual estimada del incidente.
\end{formulabox}

\subsection{1.4. Bloques de Comandos y Código Técnico}
\begin{lstlisting}[style=monocodestyle, language=bash]
# Verificacion de integridad y permisos en Linux
chmod 600 /etc/security/audit.conf
chown root:root /etc/security/audit.conf
auditctl -l -a always,exit -F arch=b64 -S open,openat -F success=0
\end{lstlisting}

\end{document}
